\section{The actual verification frontier}\label{a-verify-frontier}

This companion records what has been proved, what has been computed, and what remains to connect the adic program to the formal physics archive. The finite construction uses ordinary complex amplitudes on finite address sets. Its exact statements concern unit coordinates, normalized embeddings, packet projectors, compatible Hamiltonians, Schur complements, and analytic residue lanes. The infinite construction has additional measure, domain and convergence hypotheses.

The adic source correspondence is limited to related modules in the preserved Block 49 archive. No Lean invocation was made for that adic audit, and its archived source was not changed. The separately checked free-Dirac generator does not certify the adic spectral program. Eight related modules supply normalized cell reconstruction, measure-preserving address coding, actual dyadic projections, aggregation composition, and a complete-metric refinement bridge. Their thirty-one selected declaration locations are recorded in the source map. They do not constitute a formalization of the complete adic spectral program.

The verification title is therefore deliberate. The monograph contains ordinary mathematical proofs, the atlas contains fresh bounded numerical checks, and this companion adds fourteen exact finite assertions. The inherited computational ledgers are retained separately. A successful identity check does not authenticate a global infinite-depth claim, and an archive hash does not authenticate a new kernel execution.

\subsection{The distinction that controls the collection}

A normalized embedding satisfies \(I_h^*I_h=1\). A compatible self-adjoint Hamiltonian also satisfies
\[
 H_{h+1}I_h=I_hH_h.
\]
Only after both properties hold does the finite functional-calculus conclusion follow. A unitary fine walk alone does not supply this intertwining. In particular, extending the depth of the archived Grover product graph changes degrees and fails under the natural inclusion of its old arcs. Section~\ref{a-verify-grover} gives an exact witness. This failure concerns that inclusion, rather than unitarity at either fixed depth.

The portable crosswalk is \path{editorial/adic/verification_source_map.json}. It pins 116 selected archive members, including source chapters, executable programs, data, formal modules and receipt logs. Every member is identified by its original ZIP path and SHA-256 digest. These are evidence locators, not 116 independent mathematical results.

\clearpage
\section{Claims and stable source locations}\label{a-verify-claims}

The technical supplement groups this material under C15 and C16. Its printed chapter ranges are stale relative to the current source package. The filenames below are stable locators. A chapter title or a compiled position must not replace an exact member path when checking the evidence.

\begin{longtable}{@{}L{.07\textwidth}L{.38\textwidth}L{.45\textwidth}@{}}
\toprule
ID & Source family & Exact scope of the current result\\
\midrule\endhead
C15 & Unit registers, chapters 29--31 & Finite unit-coordinate bijections, CRT permutation unitaries and compatible multi-prime Hamiltonians. Complex Born amplitudes.\\
C15 & Normalization and product walks, chapters 32--33 & Counting-norm normalization and unitarity at each fixed product graph. Depth compatibility requires its own proof.\\
C15 & Radial and polyphase reduction, chapters 34--35 & Finite Jacobi and block reductions. Schur formulas hold where the eliminated resolvent exists.\\
C16 & Infinite Haar limit, chapter 36 & Dense cylinder union, strong projection convergence, and self-adjoint diagonal hierarchy with its actual domain.\\
C16 & Prime decimation, chapter 37 & Exact finite Schur composition on common regular charts, with a polynomial trace-coordinate composition law.\\
C16 & Companion packet lanes, chapter 41 & Entire residue-lane identities and bounded operator functions. Unbounded operator extensions need separate domain control.\\
\bottomrule
\end{longtable}

Here ``chapter 29'' means the source filename beginning \path{29_}, rather than a claim about the current compiled chapter number. The full names are recorded in the machine-readable crosswalk. C15's executable path is precisely V1E38--V1E44, seven programs. C16 selects V1E45, V1E46, V1E49 and V1E50, four programs. Together they contribute eleven retained ledgers.

The monograph also makes the complex-coupling convention explicit: a reverse transition must carry the complex conjugate coefficient for a Hermitian Hamiltonian. A real-coefficient example cannot verify an unrestricted complex-coefficient statement. Likewise a regular radial branch relation requires its energy to lie outside the eliminated tail spectrum.



\clearpage
\section{What the inherited ledgers actually record}\label{a-verify-ledgers}

The eleven selected CSV ledgers contain 143 rows labelled PASS. The count is a historical source inventory. No original program was rerun to produce these rows for the documented audit. Exact copies are provided in \path{evidence/adic/inherited_ledgers/}, and the source map records their archive identities.

\begin{longtable}{@{}L{.12\textwidth}L{.45\textwidth}L{.1\textwidth}L{.23\textwidth}@{}}
\toprule
Code & Result family & Rows & Largest recorded residual\\
\midrule\endhead
38 & Unit coordinates and CRT & 8 & \(7.45\times10^{-15}\)\\
39 & Finite Haar refinement & 8 & \(1.10\times10^{-14}\)\\
40 & Finite-adelic dynamics & 9 & \(9.93\times10^{-9}\)\\
41 & Normalization weights & 26 & \(1.89\times10^{-15}\)\\
42 & Product graph walk & 7 & \(1.31\times10^{-15}\)\\
43 & Radial self-energy & 9 & \(2.07\times10^{-14}\)\\
44 & Polyphase refinement & 25 & \(1.27\times10^{-14}\)\\
45 & Infinite-limit surrogates & 12 & \(3.71\times10^{-14}\)\\
46 & Prime decimation & 14 & \(1.02\times10^{-12}\)\\
49 & Hypergeometric packets & 17 & \(3.91\times10^{-16}\)\\
50 & Transfer trace coordinates & 8 & \(7.77\times10^{-12}\)\\
\bottomrule
\end{longtable}

The residuals are heterogeneous. Some are matrix norms, others are coefficient differences, fitted-order discrepancies, or zero entries attached to an analytic explanation. Code 40's largest value is an order-fit discrepancy, rather than a failed exact commutator. Its PASS labels must be read with each row's method and note. Comparing all eleven maxima to one universal tolerance would have no mathematical meaning.

Code 38 checks a complete local coordinate table but samples the multiplication law. Code 43 evaluates the visible branch relation only for eigenvectors with visible weight greater than \(0.08\). Code 45 evaluates finite cylinders, not an infinite matrix on its domain. Those restrictions remain part of the retained evidence even when their CSV statement is shorter.

\clearpage
\section{Executable routes and their limits}\label{a-verify-executables}

The selected original programs are executable Python scripts, rather than the cached verifier wrappers used elsewhere in the gravity source. Their top-level statements construct examples, write tables and render figures. Running one imports its numerical and plotting dependencies and may perform substantial matrix work. The accompanying preparation record did not execute them as a group.

For an original reproduction, first extract the unmodified Volume I source archive into a working directory. The program resolves its root from \path{__file__}, then writes to that root's \path{data/} and \path{figures/} directories. An example command, run from the extracted source root, is
\begin{lstlisting}
python code/39_adic_haar_quantum_refinement.py
\end{lstlisting}
To reproduce the finite-cylinder experiment rather than claim an infinite execution, use
\begin{lstlisting}
python code/45_infinite_adic_haar_quantum_limit.py
\end{lstlisting}
These are reproduction routes, not commands already run for the documented audit.

Several scripts build rows with a literal PASS string and apply a threshold only after exporting the CSV. Code 44 exits if the largest residual exceeds \(2\times10^{-10}\); code 45 uses the same aggregate bound, and code 46 uses \(2\times10^{-9}\). Code 49's row builder marks any finite residual PASS, while its final check requires seventeen rows and a maximum residual below \(5\times10^{-10}\). Code 50 checks eight rows and the same bound. A prematurely written CSV can therefore retain PASS labels even if a later final guard fails.

Code 40 ends by printing its table and has no final aggregate rejection guard. Its exact algebra, perturbative bounds, finite exponential calculations and fitted orders need separate interpretations. A zero process exit alone would not establish every printed theorem.

The portable package preserves the original source bytes and outputs, while the new exact checker reports each implemented assertion explicitly. It exits nonzero if an assertion fails. Dependency versions and tolerance choices still belong to a reproduction record. The old archive does not supply a fully pinned environment for every original numerical script.

\clearpage
\section{Formal normalized cell reconstruction}\label{a-verify-formal-cells}

Block 35's \path{Relativistic3DL2CellReconstruction.lean} builds an actual map from a finite complex Hilbert state into \(L^2\) using disjoint measurable cells of volume one. The declaration \path{UnitVolumePartition.stepIsometry} is a linear isometry, rather than an abstract assumption that interpolation preserves norms. Its adjoint averaging satisfies \path{averaging_step}: interpolation followed by averaging returns the original finite state.

Block 36 replaces unit cells by cells of a positive common volume \(v\). Its normalization is
\[
 c_v=v^{-1/2},\qquad c_v^2v=1.
\]
The exact source theorem is \path{EqualVolumePartition.normalization_sq_mul_volume}. The declarations \path{inner_normalizedCellIndicator} and \path{stepIsometry} then establish the normalized inner product and the finite-to-continuum isometry.

This gives a real formal explanation of the normalization used in an equal-child adic refinement. If a parent is split into \(p\) children, counting-norm coefficients must be divided by \(\sqrt p\). On the continuum side, normalized indicators acquire the reciprocal square root of cell measure. The two descriptions agree after the appropriate cell-volume identification. The current formal module instantiates an octree, whose cell volumes are \(8^{-n}\) and side lengths \(2^{-n}\), and proves
\[
 8^{-n}=(2^{-n})^3.
\]

The source does not define \(\mathbb Z_p\), its unit group, or a prime-indexed Haar character basis in these selected modules. Generalizing the reusable normalization theorem to an actual \(p\)-adic measure model still requires the relevant cells, measure and disjointness facts. A similarity between \(8\)-child refinement and \(p\)-child refinement is useful guidance, not a source declaration that identifies the spaces.

The density theorem \path{Octree.finitePrefixSubmodule_dense} proves finite-prefix approximation in the actual octree \(L^2\) model. Its conclusion concerns approximation of states. It does not assert compatibility of an arbitrary hierarchical Hamiltonian with the projections.



\clearpage
\section{Formal address coding and measure}\label{a-verify-formal-coding}

Blocks 37 and 38 make the dyadic address-to-space bridge concrete. An octant address is converted into three binary streams, and binary expansion sends those streams into the unit cube. The source proves continuity and surjectivity of the coding map. It also proves a negative fact that must survive any interpretation: binary coding is not pointwise injective.

The theorem \path{binaryPoint_not_injective} uses the two familiar expansions of a dyadic rational. Measure theory handles this ambiguity by identifying functions almost everywhere. The later measure-preserving theorem is therefore stronger and more suitable for \(L^2\) than an incorrect claim of a pointwise bijection.

In \path{Relativistic3DBinaryLebesgue.lean}, the binary distribution argument proves its actual cumulative distribution and then measure preservation. The selected declarations include \path{binaryCDF_eq}, \path{binaryPoint_measurePreserving}, and \path{cubeCode_measurePreserving}. Pullback consequently preserves the \(L^2\) norm in the relevant measure model.

The correspondence is between a symbolic dyadic address space and Euclidean volume on the cube. The adic monograph's finite unit-coordinate CRT permutation is a different construction. It is a literal finite basis bijection, whereas binary Euclidean coding has pointwise ambiguity and is used through a measure-preserving pullback. Neither construction should inherit the other's hypotheses by analogy.

A \(p\)-adic physical interpretation would require another explicit bridge: an actual \(p\)-adic address measure, the intended observable algebra, and any proposed spacetime realization. The current monograph establishes its internal complex Hilbert dynamics without assuming that such a realization has been supplied.



\clearpage
\section{Formal projections and their topology}\label{a-verify-formal-projections}

Block 39 defines the adjoint of Euclidean-to-octree pullback and uses it to reconstruct actual Euclidean approximation spaces. The theorem \path{cubeAdjoint_pullback} gives the left-inverse property. The module defines orthogonal projections onto closed levels, proves contraction, proves nesting, and proves strong convergence to the identity.

For levels \(a\le b\), \path{euclideanProjection_of_le} has the form
\[
 P_bP_af=P_af.
\]
The theorem \path{euclideanProjection_tendsto} states
\[
 \|P_nf-f\|\longrightarrow0
 \qquad\text{for each fixed }f\in L^2.
\]
Block 40 then identifies those levels with the closed span of actual dyadic-box indicators. The theorem \path{prefixCell_ae_eq_preimage_dyadicBox} identifies prefix cells and box preimages almost everywhere. It avoids an invalid pointwise assertion at the boundary of binary expansions.

The declaration \path{euclideanLevel_eq_euclideanDyadicLevel} is a subspace identity. The theorem \path{euclideanDyadicProjection_tendsto} gives strong convergence for the box realization. Thus the formal bridge includes actual normalized interpolation, actual Euclidean cells, and convergence of the actual projections.

Strong convergence is the topology needed for state reconstruction. It is not operator-norm convergence. If an infinite-dimensional space has a nonzero detail vector above every cutoff, then
\[
 \|1-P_n\|=1
 \qquad\text{for every }n.
\]
That observation is fully compatible with \(P_nf\to f\) for each fixed \(f\). A bound on a chosen regularity class is a third statement and needs its own weighted norm assumption.

The adic infinite-limit proof uses the analogous cylinder-filtration mechanism in its explicitly defined space. The selected Lean modules establish the dyadic Euclidean mechanism. They do not yet instantiate the prime-indexed Hamiltonian, its resolvent, or the companion residue-lane functional calculus.

\clearpage
\section{The formal exact-tower interface}\label{a-verify-formal-tower}

The module \path{LocalDiscreteGlobalContinuum.lean} defines an exact refinement tower through its actual required data: state types, embeddings, level evolutions and an adjacent intertwining hypothesis. Its declaration \path{evolve_iterEmbed} propagates
\[
 U_{n+1}I_n=I_nU_n
\]
through any finite number of embeddings. It does not construct the embeddings or prove adjacent compatibility for every concrete walk.

The same module defines a \path{SummableDefect}. A reconstruction sequence \(f_n\) in a complete metric space must satisfy
\[
 d(f_n,f_{n+1})\le b_n,
 \qquad \sum_{n=0}^{\infty}b_n<\infty.
\]
The declarations \path{exists_limit} and \path{dist_limit_le_tail} then prove convergence and the quantitative bound
\[
 d(f_n,f_\infty)
 \le\sum_{j=0}^{\infty}b_{n+j}.
\]
The declared complete target and summable sequence are load-bearing assumptions. A small finite residual does not manufacture a summable defect sequence.

The theorem \path{globalEvolution_error_bound} applies the reconstruction mechanism to the coherently evolved tower. For the adic model, a concrete use would need the specified prime-indexed spaces, normalized embeddings, compatible evolution and a reconstruction satisfying the defect requirement. A diagonal hierarchy that stabilizes old Haar weights gives one ordinary mathematical route. A depth-dependent graph walk with changing local coins requires another route.

\path{RefinementTower.lean} supplies two further exact building blocks: \path{push_comp} proves composition of finite aggregation, while \path{aggregateCurrent_comp} proves composition of the matrix current transfer \(F\mapsto RFR^T\). These aggregation identities should not be renamed as Schur-complement transitivity. Aggregation sums variables; Schur elimination returns the effect of an inverted hidden block. The monograph proves the latter separately on its regular spectral chart.

\clearpage
\section{Inherited formal receipt identity}\label{a-verify-receipts}

The accompanying record compared selected current archive bytes against retained source hashes. The seven selected modules already mapped in the gauge companion match their inherited identities. The eighth, \path{RefinementTower.lean}, is pinned to the current archive and read directly; this companion does not add a new historical receipt comparison for it.

The narrowed Block 35--40 receipt audit contains 37 matching log hashes, six unresolved log paths, and three available logs without an expected receipt digest. This count concerns evidence availability and identity. It is not a new compilation count.

\begin{longtable}{@{}L{.13\textwidth}L{.77\textwidth}@{}}
\toprule
Block & Retained evidence boundary\\
\midrule\endhead
35 & Source, API, axiom, build and negative-control identities are available in the selected receipt scope.\\
36 & Three cited build or source-compile paths are missing. Available API, axiom and negative-control logs remain usable as inherited identity evidence.\\
37--39 & Selected source and retained receipt-log identities match. Block 39 uses a concrete Euclidean projection module.\\
40 & Direct compile, module build and cumulative build log paths are missing. Three negative-control logs are available but unpinned by the original receipt.\\
\bottomrule
\end{longtable}

The formal archive retains Lean \(4.34.0\)-rc2 and the recorded Mathlib revision for these releases. The corresponding reports declare their trust and axiom policies. The accompanying record leaves that inherited compiler and cache boundary unchanged. Its source map does not infer freshness from the phrase ``checked archive.''

A useful future kernel check would instantiate a specific finite Haar embedding and its Hamiltonian compatibility first. It should run the actual declaration and API check with the recorded toolchain, produce fresh complete logs, and pin those logs to exact source bytes. Recompiling unrelated formal modules would not close that instantiation gap.

An unresolved log is reported as missing evidence, rather than a mathematical counterexample. Conversely, a matching empty successful compile log is an identity match to its historical receipt, rather than proof that a compile just occurred. These distinctions let the collection preserve progress without upgrading incomplete provenance.

\clearpage
\section{The fourteen fresh exact assertions}\label{a-verify-fresh}

The script \path{tools/verify_adic_identities.py} performs small deterministic checks in integers, rational numbers and exact symbolic rational functions. It does not execute any original archive program. Its final bounded run reports fourteen passed assertions and zero failed assertions. Three of the checks are named negative controls: their success means that a proposed incorrect simplification was detected.

\begin{longtable}{@{}L{.05\textwidth}L{.55\textwidth}L{.3\textwidth}@{}}
\toprule
No. & Implemented assertion & Tested range\\
\midrule\endhead
1--2 & Rational averaging algebra and detail ranks & \(p=2,3\), depth 3, at most \(27\times27\)\\
3 & Incidence Gram normalization & \(p=2,3,5\), coarse dimension 3\\
4--5 & CRT dual phase and missing-dual negative control & Moduli \((2,3),(3,5),(4,9)\), all 1557 basis pairs\\
6 & Cyclic-shift incidence intertwining & \(p=2,3,5\), coarse dimension 3\\
7--9 & Resolvent return, coupling degree and determinant factorization & One exact symbolic \(3\times3\) family\\
10 & Nested Schur transitivity & One exact rational \(4\times4\) matrix\\
11 & Grover depth-extension leakage witness & \(p=2,3\), old degree 3\\
12 & Factorial residue-lane coefficients & \(2\le N\le7\), all \(j\), \(0\le q\le8\)\\
13--14 & Reducing packet and coupled negative control & Rational dimension 4, polynomial degree 3\\
\bottomrule
\end{longtable}

The negative controls detect the incorrect CRT output permutation, nonzero Grover depth-extension leakage, and off-diagonal packet coupling. The checker reports fourteen assertion groups, not fourteen unrelated theorems.

The current result and manifest are in \path{evidence/adic/identity_checks_final/}. The version 2 manifest verifies the final script hash before and after execution, declared result hashes, exit status and effective resource caps. It was validated against the live project root. Correct interpretation still depends on the bounded contract and the displayed ordinary proofs.

\clearpage
\section{The CRT phase check and normalization}\label{a-verify-crt}

For coprime \(a,b\) and \(N=ab\), the input CRT map sends \(n\) to \((n\bmod a,n\bmod b)\). The Fourier output must instead use
\[
 k_a=kb^{-1}\pmod a,\qquad
 k_b=ka^{-1}\pmod b.
\]
Indeed the product Fourier kernel has exponent
\[
 2\pi i\left(\frac{k_an_a}{a}+\frac{k_bn_b}{b}\right).
\]
After putting over denominator \(ab\), its numerator is congruent to \(kn\) modulo \(ab\). Thus it equals the \(N\)-point kernel. The exact checker tests this congruence for all basis pairs of three selected coprime modulus pairs. It also detects that reusing the input residue map for the output fails in this family.

The test uses integer modular arithmetic, rather than floating complex exponentials. It therefore separates a phase convention mistake from a tiny numerical Fourier residual. The general congruence argument supplies the proof for arbitrary coprime moduli; enumeration supplies a bounded implementation check.

The normalization check has a similarly exact surrogate. Let \(B\) be the coarse-to-fine incidence matrix with one parent assigned to each fine address and \(p\) children per parent. Then
\[
 B^*B=p\,1,
 \qquad I=p^{-1/2}B,
 \qquad I^*I=1.
\]
The rational checker verifies \(B^*B=p1\) directly and keeps the square-root normalization outside the integer matrix. Raw replication increases squared norm by \(p\); it must not be used as a counting-norm isometry.

For a common finest depth \(H\), the conditional averaging matrices are rational:
\[
 (E_r)_{xy}=p^{-(H-r)}
 \mathbf1_{x\equiv y\ ({\rm mod}\ p^r)}.
\]
The checks give \(E_rE_s=E_{\min(r,s)}\) and the orthogonal differences \(D_r=E_r-E_{r-1}\). Their traces recover the detail ranks. These exact formulas supply independent evidence alongside the atlas's floating Haar-basis computations.

\clearpage
\section{Return terms and regular spectral charts}\label{a-verify-schur}

The exact symbolic family used by the checker is
\[
 H=\begin{pmatrix}
 1&\lambda&0\\ \lambda&2&1\\0&1&3
 \end{pmatrix}.
\]
Retain the first coordinate and eliminate the remaining two. Direct inversion of the detail block gives
\[
 \Sigma(z)=\lambda^2
 \frac{z-3}{(z-2)(z-3)-1},
 \qquad
 \bigl((z-H)^{-1}\bigr)_{00}
 =\frac1{z-1-\Sigma(z)}.
\]
The resolvent formula assumes both the detail inverse and the full inverse exist. The determinant factorization can be cleared to a polynomial identity, but a cleared identity does not make every intermediate inverse defined at a pole.

The checker obtains the self-energy independently from the block inverse, verifies the displayed compression of the full inverse, and verifies determinant factorization. It also compares direct and staged elimination of a separate rational \(4\times4\) matrix. That comparison tests finite Schur transitivity under invertibility, rather than an infinite operator limit.

The quadratic factor \(\lambda^2\) has a precise perturbative interpretation. If \(\lambda=p^{-w}\), the self-energy is of order \(p^{-2w}\) with a constant controlled by the detail resolvent. Letting \(z\) approach a detail pole can destroy a uniform bound. A fitted small-coupling slope away from the pole does not close that issue.

The packet negative control addresses a different possible shortcut. A block-diagonal rational \(4\times4\) Hamiltonian preserves its rank-two reducing packet under a cubic polynomial function. Adding a symmetric visible-to-hidden entry produces nonzero leakage under the same polynomial. Residue-lane identities alone do not force a spatial or packet projector to reduce an arbitrary Hamiltonian. The commutation or intertwining hypothesis remains essential.

\clearpage
\section{A concrete depth-compatibility failure}\label{a-verify-grover}

The product graph in the source combines a cycle with a rooted regular tree and uses the Grover coin at each vertex. At a fixed depth, the coin and arc-reversal shift are unitary. Extending the tree depth changes the degree of an old leaf. This prevents the natural inclusion of old arcs from being an exact refinement intertwiner.

At an old leaf away from exceptional cycle sizes, the old degree is \(d=3\): two cycle arcs and one arc to its parent. After extension, \(p\) child arcs are added. The new coin is
\[
 G_{d+p}=\frac{2}{d+p}\mathbf1\mathbf1^*-1.
\]
Apply it to one old outgoing arc basis vector. Every newly added outgoing child arc receives amplitude \(2/(d+p)\). Therefore the squared norm in the new-arc sector is
\[
 \|QG_{d+p}e_j\|^2
 =\frac{4p}{(d+p)^2}>0.
\]
The shift only permutes arcs and preserves this norm. For the checked cases,
\[
 p=2:\quad\|QG_5e_j\|^2=\frac8{25},
 \qquad
 p=3:\quad\|QG_6e_j\|^2=\frac13.
\]
The coarse evolution followed by old-arc inclusion cannot contain these new arcs, so the intertwining defect is nonzero.

This is an exact finite counterexample to one proposed embedding route. It does not invalidate fixed-depth probability conservation, the graph-distance propagation cone, or the hierarchical Haar Hamiltonian whose old sectors are explicitly stabilized. Those are different results with different operators.

A compatible graph tower could use different embedding data, boundary treatment, local coins, or quantitative approximate compatibility. Any such route needs a construction and proof. The complete-metric exact-tower theorem cannot accept a failed compatibility hypothesis. The atlas's fixed-depth product-walk example should consequently be read as a unitary finite graph illustration, with no automatic continuum trajectory.

\clearpage
\section{Infinite depth: the domain and topology gate}\label{a-verify-infinite}

In the monograph's cylinder realization, let \(P_h\) project onto the level-\(h\) subspace, and let \(D_r\) be its orthogonal detail differences. Density of the cylinder union gives \(P_hf\to f\) in \(L^2\). A hierarchical Hamiltonian with real detail weights \(\lambda_r\) is an unbounded operator when those weights grow. Its definition includes
\[
 \Dom(H)=
 \left\{f:\sum_r\lambda_r^2\|D_rf\|^2<\infty\right\},
 \qquad Hf=\sum_r\lambda_rD_rf.
\]
A state can belong to \(L^2\) and fail this domain condition. Strong convergence of cylinder projections alone does not imply convergence of \(HP_hf\).

For example, choose one unit vector \(e_r\) from each nonzero detail sector and coefficients \(1/(r+1)\). Then \(\sum_r(r+1)^{-2}<\infty\), so the state is square integrable. With weights \(\lambda_r=r+1\), its graph-norm sum contains infinitely many terms equal to one. It lies outside \(\Dom(H)\). The obstruction is an ordinary infinite-series argument, not a numerical cutoff test.

For a nonreal spectral parameter, a self-adjoint diagonal hierarchy has a bounded resolvent. Finite cutoff resolvents converge strongly on fixed states when the compatible hierarchy is used. On the omitted sectors the resolvent norm is controlled by
\[
 \sup_{r>h}|z-\lambda_r|^{-1}.
\]
If the weights tend to infinity in absolute value, this tail tends to zero. If the weights stay bounded, that particular operator-norm tail conclusion need not hold. A finite-state tail estimate and an operator-norm tail estimate are separate statements.

The original code 45 studies finite cylinders and prescribed coefficients. Its small residuals are useful checks of a finite surrogate. The uniform infinite conclusion comes from the monograph's orthogonal-series proof with explicit domain, rather than from the largest tested depth. The selected Lean bridge formalizes related dyadic projections, not this prime-indexed spectral operator.

\clearpage
\section{Companion lanes and bounded operator functions}\label{a-verify-lanes}

For \(N\ge2\) and \(0\le j<N\), the exponential residue lane is
\[
 E_{N,j}(z)=\sum_{q=0}^{\infty}
 \frac{z^{Nq+j}}{(Nq+j)!}.
\]
The source's root-of-unity filter separates these lanes exactly, and their sum reconstructs \(e^z\). The factorial coefficient factorization is
\[
 (Nq+j)!=j!\,N^{Nq}q!
 \prod_{\alpha\in\mathcal A_{N,j}}(\alpha)_q,
\]
where \(\mathcal A_{N,j}\) contains \((j+r)/N\), \(1\le r\le N\), with its single parameter equal to one removed. This is the coefficient identity behind the generalized hypergeometric expression.

The exact checker evaluates all 243 selected coefficient instances: \(2\le N\le7\), every residue \(j\), and \(0\le q\le8\). It uses rational products, rather than floating special-function evaluations. The monograph supplies the general product argument. The atlas separately compares high-precision series, root sums and hypergeometric values for selected inputs.

For a bounded operator \(A\), the entire series defining \(E_{N,j}(A)\) converges in operator norm. The exponential reconstruction then holds in the bounded-operator algebra. If a projector \(P\) commutes with \(A\), it commutes with every norm-convergent power series in \(A\). This is the actual reason an established reducing packet survives the lane operation.

For an unbounded self-adjoint operator, \(e^{itH}\) is defined by spectral calculus for real time. A root filter containing factors such as \(e^{\omega_N^\ell zH}\) can include exponentials with unbounded real growth. Their domains may differ, even when a sum has a formally attractive scalar identity. The bounded result must not be transferred to all states of an unbounded hierarchy without a common-domain or spectral-integrability theorem.

The finite cyclic companion matrix realizes an analytic residue system. Its word ``companion'' does not identify it with a physical quantum subsystem or an additional Hilbert-space degree of freedom by itself.

\clearpage
\section{Reproduction and the next formal bridge}\label{a-verify-reproduction}

The new exact checker can be run from the bundle root with
\begin{lstlisting}
python tools/verify_adic_identities.py --output \
  evidence/adic/identity_reproduction/results.json
\end{lstlisting}
This command computes the same finite assertions and writes a new result file. A result file alone is not a fresh provenance manifest. For a recorded run, use the computation-audit runner with \path{evidence/adic/identity_contract.json}, pin \path{tools/verify_adic_identities.py}, declare the new result path, and choose a fresh output directory. The retained final manifest records the actual argv, software versions, caps, logs and SHA-256 digests.

The final run used a 45-second wall limit, a 35-second CPU cap, a 1 GiB address-space cap, one effective CPU core, conventional one-thread numerical-library settings, and a 1 MiB log limit. It completed successfully and its version 2 manifest validated against the project. Conventional thread settings can be ignored by libraries; the recorded risk remains explicit.

An initial run failed while serializing SymPy integer metadata to JSON. It produced no scientific result file. The serializer was corrected to write ordinary integers. That failure log and its exact input snapshot are retained. A subsequent successful run was superseded after the self-energy check was strengthened to compare an independently computed detail-block inverse. The final live source is pinned by \path{identity_checks_final/manifest.json}; \path{identity_run_history.json} explains all three attempts.

The next formal milestone is specific: define one finite prime-address space, its counting inner product, normalized child embedding and detail projectors; prove the projector algebra and the compatible hierarchical Hamiltonian; then prove resolvent or unitary intertwining. A later phase can formalize a prime-indexed measure realization, the dense cylinder union, the unbounded operator domain, and the spectral limit. The fixed-depth Grover graph needs a separately chosen refinement route.

Four later physical questions remain outside that milestone: identifying a spacetime realization, deriving a coupled fiber/backreaction model, supplying continuum field equations for that model, and testing physical predictions. The accompanying record provides a carefully stated mathematical substrate and reusable formal interfaces. None of its evidence counts closes those additional implications.
